% ---
% school: "college"
% level: "8eme"
% section: null
% chapter: "العمليات والحساب على الأعداد الكسرية"
% source: "Axiomaths"
% title: "حسابات ذكية"
% country: "Tunisie"
% year: 2026
% professor: "Axiomaths"
% difficulty: "Moyen"
% points: null
% tags: ["fractions", "calcul-astucieux", "factorisation", "produit-telescopique"]
% ---
% Exercice 9/18 de la série 8ème « العمليات والحساب على الأعداد الكسرية » (statement + \axsol + solution). Compiler avec XeLaTeX + axiomath.sty (+ axiomaths-logo.png).
\documentclass[12pt,a4paper]{article}
\usepackage[level=8b,color=15803D]{axiomath}
\axsetup{title={العمليات والحساب على الأعداد الكسرية}, corrige}

\begin{document}

\begin{corr}[title={حسابات ذكية}, stars=3, source={Axiomaths}]
احسب بأيسر طريقة (بالتفكيك إلى جذاء أو بالاختزال) :
\[A=\dfrac{7}{9}\times\dfrac{5}{11}+\dfrac{7}{9}\times\dfrac{6}{11}\qquad B=\dfrac{13}{8}\times\dfrac{3}{5}-\dfrac{5}{8}\times\dfrac{3}{5}\]
\[C=\dfrac{2}{3}\times\dfrac{3}{4}\times\dfrac{4}{5}\times\dfrac{5}{6}\times\dfrac{6}{7}\qquad D=\dfrac{-4}{15}\times\dfrac{7}{3}+\dfrac{-4}{15}\times\dfrac{2}{3}\]
\[E=\left(1-\dfrac{1}{2}\right)\times\left(1-\dfrac{1}{3}\right)\times\left(1-\dfrac{1}{4}\right)\times\cdots\times\left(1-\dfrac{1}{10}\right)\]
\axsol
\begin{prop}[التفكيك إلى جذاء]
\ib{$k\times a+k\times b=k\times(a+b)$} \qquad و \qquad \ib{$k\times a-k\times b=k\times(a-b)$}
\end{prop}
\begin{enumerate}[itemsep=14pt]
  \item[$A$ :] العامل المشترك هو $\dfrac{7}{9}$ :
  \[A=\dfrac{7}{9}\times\left(\dfrac{5}{11}+\dfrac{6}{11}\right)=\dfrac{7}{9}\times\dfrac{11}{11}=\dfrac{7}{9}\times1=\dfrac{7}{9}\]
  \item[$B$ :] العامل المشترك هو $\dfrac{3}{5}$ :
  \[B=\dfrac{3}{5}\times\left(\dfrac{13}{8}-\dfrac{5}{8}\right)=\dfrac{3}{5}\times\dfrac{8}{8}=\dfrac{3}{5}\]
  \item[$C$ :] نكتب الجذاء على شكل كسر واحد :
  \[C=\dfrac{2\times3\times4\times5\times6}{3\times4\times5\times6\times7}\]
  نشطب $3$ و $4$ و $5$ و $6$ من البسط والمقام :
  \[C=\dfrac{2}{7}\]
  \item[$D$ :] العامل المشترك هو $\dfrac{-4}{15}$ :
  \[D=\dfrac{-4}{15}\times\left(\dfrac{7}{3}+\dfrac{2}{3}\right)=\dfrac{-4}{15}\times\dfrac{9}{3}=\dfrac{-4}{15}\times3=-\dfrac{12}{15}=-\dfrac{4}{5}\]
  \item[$E$ :] كل عامل هو من الشكل $1-\dfrac{1}{k}=\dfrac{k-1}{k}$ :
  \[E=\dfrac{1}{2}\times\dfrac{2}{3}\times\dfrac{3}{4}\times\cdots\times\dfrac{9}{10}\]
  بالاختزال المتتالي لا يبقى إلا البسط الأول والمقام الأخير :
  \[E=\dfrac{1}{10}\]
\end{enumerate}
\textbf{الفكرة :} قبل أن نحسب، نبحث عن عامل مشترك أو عن حدود تُختزل فيما بينها.
\end{corr}

\end{document}
